\chapter{工程证据索引与提交清单}
\section{路径、版本与引用约定}
除本目录的 \path{problem.md}、\path{problem.pdf} 和报告文件外，
本附录及参考资料中的工程路径均相对
\path{/home/yzy/code/ECSDC_FPGA_2026/Pgl50h/}。
为便于核对，各章“工程依据”也列出关键源文件；引用工程摘要时，应尽量进一步核对原始报告和日志。

\begin{longtable}{p{0.10\textwidth}p{0.82\textwidth}}
\caption{报告各部分的主要证据}\label{tab:evidence}\\
\toprule 章节 & 源码、记录和用途 \\
\midrule\endfirsthead
\toprule 章节 & 源码、记录和用途 \\
\midrule\endhead
第 1 章 & 本报告目录的 \path{problem.md}；工程 \path{README.md}：需求与完成边界。 \\
第 2 章 & \path{EmbFPGA.pds}、\path{src/top/hdmi_ddr3_top.v}：器件、工程入口、实际默认配置。 \\
第 3 章 & \path{calibration/camera_20260930.json}；\path{src/video/pose/pose_inverse_project_q14.v}：标定、坐标、旋转、畸变及定点实现。 \\
第 4 章 & \path{src/video/pose/pose_renderer_tiled.v}、\path{fram_buf_pose.v}、\path{pose_ddr_arbiter.v}（同目录）；\path{src/constraints/hdmi_ddr3_top.fdc}：缓存、所有权与约束。 \\
第 4 章 & \path{verification/pose_tiled_20261010/source_bounds/README.md}：历史软件模型的源缓存容量边界及适用范围。 \\
第 5 章 & \path{src/imu/imu_protocol_50m.v}；\path{src/control/}；\path{src/video/preprocess/}；\path{report/pc_control_protocol.md}：姿态、控制和增强。 \\
第 6 章 & \path{tests/}；\path{verification/pds_pc_control_write_ram_20261022/simulation/result.log}；同一归档下的 \path{pc_control/}、\path{main_rtl/}：原始测试日志。 \\
第 7 章 & \path{bitstream/pc_control_20261022_identity.json}；\path{verification/20261022_pc_control_sbit.md}：当前配置、资源、时序及吞吐结果。 \\
第 7 章 & \path{verification/pds_pc_control_write_ram_20261022/prj_tasks/syn_1/synthesize/hdmi_ddr3_top.snr}：综合资源报告。 \\
第 7 章 & 同一归档下 \path{prj_tasks/pnr_1/report_timing/hdmi_ddr3_top.rtr}、\path{pose_fifo_cdc.json}：最终时序及 Gray 偏斜审计。 \\
第 7 章 & \path{verification/20261010_pose_tiled_30fps.md}：历史 $128\times64$ 候选；不可替代当前配置证据。 \\
\bottomrule
\end{longtable}

\section{当前位流身份}
候选路径为 \path{bitstream/EmbFPGA_1920_pc_control_20261022.sbit}，
大小 2,101,696 bytes，器件 PGL50H-FBG484-6。身份文件为
\path{bitstream/pc_control_20261022_identity.json}，SHA256 为
\begin{center}\small
\path{492e1c3c29606e4e29eb47e5f95704fb1892cc3cc0a970f555cbb3e9bd320b4c}
\end{center}
模板整理时已对实际候选文件计算并核对上述哈希。
当前归档标记为尚未烧录、尚未板测。
归档 \path{source_status.json} 记录构建当时输入快照与源码一致；
后续修改应重新核对源码身份，不能将历史状态直接作为新版本审计。

\section{最终提交材料登记表}
\begin{longtable}{p{0.29\textwidth}p{0.63\textwidth}}
\toprule 材料 & 位置、版本或待补内容 \\
\midrule\endfirsthead
\toprule 材料 & 位置、版本或待补内容 \\
\midrule\endhead
报告正文及可编辑源文件 & 本目录 \path{main.tex}、各章节；\fillin{最终版本号及日期} \\
最终 FPGA 主工程 & \path{EmbFPGA.pds}；\fillin{提交源码身份} \\
目标位流与校验记录 & \fillin{最终 sbit、SHA256、对应配置和下载记录} \\
真实标定数据 & \fillin{标定图集、标定板信息、重投影误差、参数导出脚本} \\
仿真与实现报告 & \fillin{最终配置对应的测试、资源、时序及 CDC 日志} \\
上位机及外部控制固件 & \fillin{软件位置、版本、依赖、串口与硬件接线说明} \\
实板测试原始资料 & \fillin{帧号、延迟、功耗、稳定性及姿态响应记录} \\
图像质量分析资料 & \fillin{原图、参考、输出、ROI、直线度及 MTF/SFR 曲线} \\
现场演示资料 & \fillin{连接图、操作脚本、图卡清单、演示视频及分工} \\
\bottomrule
\end{longtable}

\section{测试记录的建议字段}
\begin{center}
\small
\begin{tabularx}{\textwidth}{p{3cm}Y}
\toprule
字段 & 记录内容 \\
\midrule
test\_id / date & \fillin{唯一编号、日期、测试人员} \\
source / bitstream & \fillin{源码版本、位流哈希和工具版本} \\
camera / calibration & \fillin{相机、曝光、焦距、标定文件版本} \\
pose / mode & \fillin{yaw/pitch、机械角度、功能掩码及 Gamma 曲线} \\
stimulus / measurement & \fillin{图卡、输入方式、测量方法及设备精度} \\
result / evidence & \fillin{原始数据、统计方法、判据和图像/日志位置} \\
\bottomrule
\end{tabularx}
\end{center}
